\paragraph{积分沿真实路径，不能只代入首末位置}
上式的积分还包含路径信息。若已知运动参数 $\vect r(t)$，一般情形应写为
\[
 W=\int_{t_A}^{t_B}\vect F(\vect r(t),t)\cdot\dot{\vect r}(t)\dd t.
\]
图~\ref{quality:08:work}中 $\dd\vect r$ 沿轨迹切线，故只有 $F\cos\theta$ 进入功；法向分力即使改变运动方向，也不在这一微小位移上做功。若使用 $\int F_x(x)\dd x$，必须确认其他分量没有贡献，且已把沿该段路径的 $F_x$ 正确表示为 $x$ 的函数。
保守力的功可进一步化为势能差，但这需要力场具有相应的路径无关性，不能仅因题中给出了起点与终点就省去路径。例如滑动摩擦功是 $-\int\mu N\dd s$，积分用路程元 $\dd s\ge0$；往返运动的净位移可以为零，摩擦耗散却不为零。
\begin{qualityfigure}\begin{tikzpicture}[>=Latex]
 \draw[thick] (0,0) node[left] {$A$} .. controls (1,1.7) and (1.8,.55) .. (2.5,.55)
 .. controls (3.2,.55) and (4,.4) .. (5,1.4) node[right] {$B$};
 \fill (2.5,.55) circle (1.5pt);
 \draw[vector] (2.5,.55)--(4.1,.55) node[right] {$\dd\vect r$};
 \draw[vector,draw=TextGray] (2.5,.55)--(3.6,1.75) node[above] {$\vect F$};
 \node at (3.1,.95) {$\theta$};
 \node[below] at (2.5,-.5) {$\delta W=F\cos\theta\,\dd s$};
\end{tikzpicture}
\figcap{路径积分在每个位置取力与局部切向位移的点乘，箭头长度只作局部方向示意。}\label{quality:08:work}
\end{qualityfigure}
